%% 
%% Copyright 2007-2025 Elsevier Ltd
%% 
%% This file is part of the 'Elsarticle Bundle'.
%% ---------------------------------------------
%% 
%% It may be distributed under the conditions of the LaTeX Project Public
%% License, either version 1.3 of this license or (at your option) any
%% later version.  The latest version of this license is in
%%    http://www.latex-project.org/lppl.txt
%% and version 1.3 or later is part of all distributions of LaTeX
%% version 1999/12/01 or later.
%% 
%% The list of all files belonging to the 'Elsarticle Bundle' is
%% given in the file `manifest.txt'.
%% 
%% Template article for Elsevier's document class `elsarticle'
%% with harvard style bibliographic references

\documentclass[preprint,12pt]{elsarticle}

%% Use the option review to obtain double line spacing
%% \documentclass[preprint,review,12pt]{elsarticle}

%% Use the options 1p,twocolumn; 3p; 3p,twocolumn; 5p; or 5p,twocolumn
%% for a journal layout:
%% \documentclass[final,1p,times]{elsarticle}
%% \documentclass[final,1p,times,twocolumn]{elsarticle}
%% \documentclass[final,3p,times]{elsarticle}
%% \documentclass[final,3p,times,twocolumn]{elsarticle}
%% \documentclass[final,5p,times]{elsarticle}
%% \documentclass[final,5p,times,twocolumn]{elsarticle}

%% For including figures, graphicx.sty has been loaded in
%% elsarticle.cls. If you prefer to use the old commands
%% please give \usepackage{epsfig}

%% The amssymb package provides various useful mathematical symbols
\usepackage{amssymb}
%% The amsmath package provides various useful equation environments.
\usepackage{amsmath}
%% The amsthm package provides extended theorem environments
%% \usepackage{amsthm}

%% The lineno packages adds line numbers. Start line numbering with
%% \begin{linenumbers}, end it with \end{linenumbers}. Or switch it on
%% for the whole article with \linenumbers.
\usepackage{lineno}
\usepackage{multirow}
\usepackage{url}
\usepackage{caption}

\usepackage{listings}
\usepackage{xcolor}
\usepackage{hyperref}

\definecolor{turtle-green}{RGB}{0,128,0}
\definecolor{turtle-blue}{RGB}{0,0,255}
\definecolor{turtle-red}{RGB}{160,0,0}


\definecolor{sparql-keyword}{RGB}{0,0,180}
\definecolor{sparql-variable}{RGB}{160,0,0}     
\definecolor{sparql-uri}{RGB}{0,100,0}
\definecolor{sparql-string}{RGB}{180,0,180}
\definecolor{sparql-comment}{RGB}{120,120,120}
\definecolor{sparql-prefix}{RGB}{180,100,0}


\lstdefinelanguage{Turtle}{
  morekeywords={@prefix, @base},
  sensitive=true,
  comment=[l]{\#},
  morecomment=[s]{/*}{*/},
  morestring=[b]",
  morestring=[b]',
  moredelim=[s][\color{turtle-blue}]{<}{>},
  moredelim=[s][\color{turtle-green}]{:}{\ },
  moredelim=[s][\color{turtle-red}]{she:}{\ },
}


\lstdefinelanguage{SPARQL}{
  keywords={SELECT, WHERE, FILTER, OPTIONAL, UNION, GROUP, BY, ORDER, LIMIT, OFFSET, 
    DISTINCT, REDUCED, CONSTRUCT, ASK, DESCRIBE, FROM, NAMED, GRAPH, SERVICE, 
    BIND, VALUES, MINUS, NOT, EXISTS, AS, COUNT, SUM, MIN, MAX, AVG, HAVING, 
    DATATYPE, LANG, BOUND, sameTerm, isURI, isIRI, isLiteral, isBlank, isNumeric, 
    str, lang, datatype, REGEX, REPLACE, a},
  sensitive=true,
  comment=[l]{\#},
  morecomment=[s]{/*}{*/},
  morestring=[b]",
  morestring=[b]',
  moredelim=[s][\color{sparql-uri}]{<}{>},
  moredelim=[s][\color{sparql-variable}]{?}{\ },
  moredelim=[s][\color{sparql-variable}]{\$}{\ },
  moredelim=[s][\color{sparql-prefix}]{PREFIX}{\ },
  moredelim=[s][\color{sparql-prefix}]{prefix}{\ },
}

\lstdefinestyle{TurtleStyle}{
  language=Turtle,
  basicstyle=\ttfamily\small\selectfont,
  breaklines=true,
  frame=single,
  showstringspaces=false,
  tabsize=2,
  keywordstyle=\color{black}\bfseries,
  commentstyle=\color{gray}\itshape,
  stringstyle=\color{purple},
  identifierstyle=\color{black},
  captionpos=b,
}

\lstdefinestyle{SPARQLStyle}{
  language=SPARQL,
  basicstyle=\ttfamily\small\selectfont,
  keywordstyle=\color{sparql-keyword}\bfseries,
  ndkeywordstyle=\color{sparql-prefix}\bfseries,
  commentstyle=\color{sparql-comment}\itshape,
  stringstyle=\color{sparql-string},
  identifierstyle=\color{black},
  showstringspaces=false,
  breaklines=true,
  frame=single,
  tabsize=2,
  captionpos=b,
}

\usepackage{tcolorbox}
\tcbuselibrary{listings,breakable,skins}  % Add skins library to enable enhanced

% Define Markdown language for listings
\definecolor{md-blue}{RGB}{0,0,255}
\definecolor{md-green}{RGB}{0,128,0}
\definecolor{md-gray}{RGB}{128,128,128}
\definecolor{md-darkgray}{RGB}{80,80,80}

\lstdefinelanguage{Markdown}{
  sensitive=true,
  morecomment=[l]{\#},
  morekeywords={\#\#\#, \#\#, \#, -, *, 1., >},
  morestring=[b]",
  moredelim=[s][\color{md-blue}]{`}{`},
  moredelim=[s][\color{md-blue}\bfseries]{**}{**},
  moredelim=[s][\color{md-blue}\itshape]{*}{*},
  moredelim=[s][\color{md-green}]{```}{```},
  moredelim=[l][\color{md-gray}]{>},
}

% Create a new tcolorbox environment for Markdown prompts
\newtcblisting{markdownprompt}{
  listing only,
  listing options={
    language=Markdown,
    basicstyle=\ttfamily\small,
    breaklines=true,
    showstringspaces=false,
    tabsize=2,
    keepspaces=true,
    aboveskip=0pt,        % Remove space above listing
    belowskip=0pt,        % Remove space below listing
    lineskip=-2pt,        % Reduce line spacing (adjust as needed)
    basewidth=0.5em       % Tighter character spacing
  },
  title=System Prompt in Markdown,
  colback=gray!5,
  colframe=gray!75!black,
  fonttitle=\bfseries,
  enhanced,
  sharp corners,
  boxrule=0.5pt,
  breakable,
  top=2mm,                % Reduce top padding
  bottom=2mm,             % Reduce bottom padding
  left=2mm,               % Reduce left padding
  right=2mm               % Reduce right padding
}

\journal{Journal of Web Semantics}

\begin{document}

\begin{frontmatter}

%% Title, authors and addresses

%% use the tnoteref command within \title for footnotes;
%% use the tnotetext command for theassociated footnote;
%% use the fnref command within \author or \affiliation for footnotes;
%% use the fntext command for theassociated footnote;
%% use the corref command within \author for corresponding author footnotes;
%% use the cortext command for theassociated footnote;
%% use the ead command for the email address,
%% and the form \ead[url] for the home page:
% \title{Title\tnoteref{label1}}
% \tnotetext[label1]{}
% \author{Beichen Wang\corref{cor1}\fnref{label2}}
% \ead{beichen.wang@wur.nl}
% \fntext[label2]{}
% \cortext[cor1]{}
% \affiliation{organization={Artificial Intelligence Chair Group},
%              addressline={Wageningen University \& Research},
%              city={Wageningen},
%              postcode={6708 PB},
%              country={The Netherlands}}
% \fntext[label3]{}


\title{Accelerating the construction of an ontology-compliant soil health knowledge graph with large language models} %% Article title

% use optional labels to link authors explicitly to addresses:
\author[inst1]{Beichen Wang\corref{cor1}}
\ead{beichen.wang@wur.nl}
\cortext[cor1]{Corresponding author.}
\affiliation[inst1]{organization={Artificial Intelligence Chair Group},
             addressline={Wageningen University \& Research},
             city={Wageningen},
             postcode={6708 PB},
             country={The Netherlands}}

\author[]{Luís Moreira de Sousa\fnref{Luís Moreira de Sousa}}
\ead{luis.moreira.de.sousa@tecnico.ulisboa.pt}
\affiliation[Luís Moreira de Sousa]{organization={INESC-ID, Instituto Superior Técnico},
             addressline={University of Lisbon},
             city={Lisbon},
             country={Portugal}}

\author[inst1]{Anna Fensel}
\ead{anna.fensel@wur.nl}


% \author{} %% Author name

% Author affiliation
% \affiliation{organization={},%Department and Organization
%            addressline={}, 
%            city={},
%            postcode={}, 
%            state={},
%            country={}}

%% Abstract
\begin{abstract}
%% Text of abstract
Knowledge graphs (KGs) and ontologies provide a powerful means to organize domain knowledge in a structured and machine-readable form, but constructing ontology-compliant KGs for a new domain remains a time-consuming knowledge engineering task. In soil science, this challenge is particularly relevant because soil health knowledge is fragmented across scientific reports and other textual sources. In this paper, we present an ontology-compliant soil health KG derived from domain literature, together with a semi-automated, human-in-the-loop pipeline designed to accelerate its construction. Our pipeline leverages large language models (LLMs) to generate candidate RDF triples from text, while incorporating expert review to ensure ontological compliance and factual accuracy. To assess the quality of LLM-generated triples, we compared generated graphs against human-annotated gold standards using multiple complementary metrics. The results show that while exact Triple Match $F1$ remains low (0.2506), Graph BERTScore $F1$ reaches 0.8600, indicating strong semantic alignment despite structural differences. The mean normalized Graph Edit Distance (0.4730) further indicates such structural divergence, highlighting the need for human refinement in ontology-aligned KG construction. The resulting KG integrates unstructured knowledge into 11,719 RDF triples describing 2,017 entities, including 1,785 soil-related concepts. Its fidelity was reviewed through a competency-question-based validation process involving soil scientists. We demonstrate the KG’s use both as the semantic backbone for a knowledge discovery system in soil science and as a benchmarking dataset for the text-to-KG task. The KG, supporting ontology, and the source code of the pipeline are made open-source.
\end{abstract}


%%Graphical abstract
%%\begin{graphicalabstract}
%\includegraphics{grabs}
%%\end{graphicalabstract}

%%Research highlights
% \begin{highlights}
% \item Ontology-compliant soil health knowledge graph (KG) integrating 11,719 RDF triples.
% \item Semi-automated pipeline leverages large language models (LLMs) for KG construction.
% \item LLMs accelerate KG construction by generating 53\% (GED 0.47) of structures correctly.
% \item LLMs achieve 0.86 semantic similarity but low structural match needs human revisions.
% \item KG serves as the semantic backbone for soil science knowledge discovery system.
% \end{highlights}

%% Keywords
\begin{keyword}
%% keywords here, in the form: keyword \sep keyword

%% PACS codes here, in the form: \PACS code \sep code

%% MSC codes here, in the form: \MSC code \sep code
%% or \MSC[2008] code \sep code (2000 is the default)
Knowledge graph construction \sep Large language models \sep Ontology-compliance \sep Soil health \sep Knowledge engineering
\end{keyword}
\end{frontmatter}

%% Add \usepackage{lineno} before \begin{document} and uncomment 
%% following line to enable line numbers
\linenumbers

%% main text
%%

%% Use \section commands to start a section
\section{Introduction}\label{sec1}
Knowledge graphs (KGs) based on the Resource Description Framework (RDF) have emerged as a powerful means to store, integrate, and manage knowledge in a flexible yet semantically rich manner \citep{10.1145/3447772}. Already widely accepted and used across multiple domains, such as life sciences \citep{Smith2007-xn}, cultural heritage \citep{10.1007/978-3-031-33455-9_30}, and e-commerce \citep{regino2024generating}, RDF KGs retain information in ``triples'' (subject, predicate, object), linking diverse knowledge sources into a coherent graph that supports queries.
Underpinning these KGs are ontologies \citep{10.1145/3447772}, which provide pre-defined classes and properties for domain entities and their relationships. By offering standardized, formally defined terms, ontologies are crucial for ensuring consistent knowledge representation tailored to a specific domain, improving queryability, and supporting inferencing and reasoning.

Despite these advantages, constructing and maintaining an ontology-compliant \citep{10.1007/978-3-031-43458-7_48} KG for a new domain remains a costly and labor-intensive knowledge engineering task \citep{KOUTSIANA2026100873}. Beyond extracting candidate facts from text, the process also requires selecting appropriate conceptual granularity, aligning entities and relations with formal ontologies, resolving structural inconsistencies, and validating whether the resulting graph supports the intended use cases. These challenges are particularly visible in scientific domains, where knowledge is often fragmented across literature, and where semantic precision matters for downstream reuse.

Soil science represents one such domain where knowledge organization through semantic technologies, particularly KGs, can yield benefits. Soil is the foundation of world food production; it performs as a climate regulator and is a key pillar of biodiversity. The dissemination and utilization of soil-related knowledge, such as research papers, scientific reports, and best practices, are essential for informed soil management and policymaking. However, as in many other fields that have not yet widely used the Semantic Web as their primary paradigm for knowledge representation \citep{DRURY2019487}, knowledge about soil health remains scattered, ambiguous, and controversial across diverse scientific literature, policy documents, and practical guidelines, making it challenging to access, integrate, and apply effectively. This knowledge is often contained within unstructured texts, and the lack of standardized formats and common frameworks for knowledge sharing hinders the ability to synthesize findings and apply them to address real-world soil health challenges. Furthermore, the concept of soil health lacks a universally agreed-upon definition and can vary in interpretation in the context of research versus policymaking and practice. These characteristics make soil health not only an important application domain, but also a demanding case for ontology-compliant knowledge engineering.

Recent advances in Large Language Models (LLMs), such as the Generative Pre-trained Transformer (GPT) \citep{openai2024gpt4technicalreport}, offer an opportunity to accelerate KG construction by automating the extraction of entities and relations from text \citep{Zhu2024}. However, automatically ensuring that the outputs from these LLM-based approaches align with established ontologies, an essential requirement for creating robust and interoperable KGs, remains a challenge \citep{bian2025llmempoweredknowledgegraphconstruction}. This is particularly pertinent in specialized scientific domains like soil science with its specific terminology and relationships. Therefore, it becomes feasible to design a semi-automated, human-in-the-loop workflow \citep{SHIMIZU2025100862}. Such a workflow can leverage the speed of LLMs for the initial extraction while incorporating human expertise for the validation and ontology alignment, balancing automation with accuracy.

Given these challenges, fragmented knowledge, lack of standardization, and the difficulties in creating ontology-compliant KGs automatically, there is a clear need not only for a dedicated soil health KG but also for a semi-automated KG construction pipeline that can effectively build such a resource. In this study, we present the soil health KG, and the human-in-the-loop, semi-automated methodology used to create it, primarily from the report \textit{Soil Monitoring in Europe} \citep{doi/10.2800/956606}, complemented where necessary by additional materials. The report is hereinafter referred to as the EU soil health report. This KG aims to provide a structured, semantically rich, and FAIR \citep{Wilkinson2016} representation of soil health knowledge. To achieve this, we proposed a pipeline that incorporates LLMs for extracting relevant information from the source text, followed by aligning with established ontologies, post-processing, and expert validation. This pipeline thereby supports knowledge engineering in a specialized scientific domain while reducing part of the manual effort involved in building a KG from scratch.

The main contributions of this work are as follows:
\begin{itemize}
    \item A semi-automated, human-in-the-loop pipeline for ontology-compliant KG construction that leverages LLMs as an accelerator for knowledge extraction while incorporating human oversight for quality assurance.
    \item A dedicated, expert-validated soil health KG comprising 11,719 RDF triples representing 2,017 entities, including 1,785 concepts, interlinked with established soil-related vocabularies and thesauri.
    \item A multi-perspective evaluation of LLM-generated graphs against human-annotated gold standards, showing that LLM outputs often preserve the semantic content of the source text while still requiring structural refinement.
    \item Two practical applications of this KG, which serves as a semantic backbone for a knowledge discovery system and constitutes a new benchmarking dataset for the text-to-KG research.
\end{itemize}

The remainder of this paper is organized as follows. Section~\ref{sec2} reviews related work on knowledge engineering in soil science and LLM-empowered KG construction. Section~\ref{sec3} describes the materials and methods, including the knowledge sources, ontology design, KG construction pipeline, and human-in-the-loop validation process. Section~\ref{sec4} presents the resulting soil health KG and the evaluation of the LLM-generated graphs. Section~\ref{sec5} discusses the main findings. Section~\ref{sec6} illustrates two practical applications of the KG, and Section~\ref{sec7} concludes the paper. \ref{app1} and~\ref{app2} provide, respectively, the system prompt used for KG generation and the formal definitions of the evaluation metrics.


\section{Related work} \label{sec2}
\subsection{Knowledge engineering in soil science}\label{sec2_1}
In the domain of soil science, ontologies provide structured frameworks essential for organizing data and knowledge, while KGs built upon them create powerful, machine-readable assets that enable more effective knowledge interpretation and automation. Notable examples include the Soil Food Web Ontology (SFWO) \citep{LEGUILLARME2023102360}, the Global Soil Information System Web Ontology (GloSIS) \citep{doi:10.1177/22104968251363767}, and the Ontology of Soil Properties and Processes \citep{10.1007/978-3-319-46547-0_4}.

Despite the existence of these ontologies, the creation of KGs that fully capture the broader dimensions of soil science is still in development. An example of progress in this direction is described by \citet{LEGUILLARME2023103497}, which presents a method for building a biodiversity KG by integrating data on soil organisms from multiple sources using the SFWO as its ontology. \citet{HAN2022115685} constructed a KG to support urban soil pollution management. More recently, the Soil Organic Carbon Knowledge Graph (SOCKG) was introduced to provide representation of agricultural research data, focusing on the integration of soil carbon dynamics with farming practices and environmental factors \citep{10.1007/978-3-031-97207-2_18}. AGROVOC \citep{SUBIRATSCOLL2022105965}, which is developed by the Food and Agriculture Organization (FAO), provides a multilingual, controlled vocabulary covering a wide range of concepts relevant to their areas of interest, including soil science.

Most of these existing efforts focus primarily on data collection, annotation and harmonization for specific soil science applications, rather than providing a systematic representation that captures the knowledge, concepts, and vocabularies central to soil health. This gap underscores the need for a dedicated soil health KG that goes beyond data integration to encompass the conceptual framework and domain knowledge essential for understanding and managing soil health holistically.


\subsection{LLM-empowered KG construction}\label{sec2_2}
Parallel to these ontological developments, the use of LLMs for KG construction from text has garnered growing attention \citep{10.1109/TKDE.2024.3352100}, \citep{bian2025llmempoweredknowledgegraphconstruction}. Studies such as \citep{zhang-soh-2024-extract}, \citep{regino2024generating}, \citep{doi:10.3233/SW-243719}, and \citep{carta2023iterativezeroshotllmprompting} highlight how LLMs can streamline entity and relation extraction, thus reducing reliance on manual curation. However, the raw KG output from LLMs frequently contains duplicate entities and relationships. KGGen \citep{mo2025kggenextractingknowledgegraphs} is proposed as a framework to normalize those synonymous entities and disambiguate redundant relationships by LLM prompting. \citet{ghanem:hal-04862235} then compare the performance of LLM fine-tuning and prompting on KG construction. Nevertheless, challenges persist in ensuring semantic interoperability and addressing potential biases introduced by LLM-driven pipelines. Integrating these automated outputs with existing domain ontologies is especially important to maintain coherent, reusable, and accurate KGs.

Recent advances in graph-based Retrieval-Augmented Generation (RAG) have demonstrated the potential of using LLMs to generate KGs for information retrieval and advanced question-answering (QA) \citep{han2025retrievalaugmentedgenerationgraphsgraphrag}. GraphRAG \citep{edge2025localglobalgraphrag} proposes deriving a KG from source documents and clustering closely related entities, thereby enabling it to address corpus-level summarization queries that naive RAG systems cannot handle, while HippoRAG \citep{10.5555/3737916.3739818} introduces a retrieval framework inspired by hippocampal memory indexing to enhance long-term integration of knowledge in LLMs. These systems leverage LLMs to automatically generate KGs from text corpora to improve document indexing and QA performance. However, the KGs produced by these approaches are typically ad-hoc and schemaless, constructed dynamically without adherence to established ontologies or formal knowledge representation standards. Moreover, we still lack effective means to evaluate the quality and factual accuracy of these automatically generated KGs \citep{xiang2025usegraphsragcomprehensive}. Consequently, while these graph-based RAG systems may showcase impressive performance on information retrieval and QA tasks, they have not yet been proven to provide reliable and transparent knowledge representations of the underlying corpus.

\section{Materials and methods} \label{sec3}
\begin{figure}
\includegraphics[width=\textwidth]{text2KG.pdf}
\caption{Knowledge graph construction pipeline. The blue modules indicate manual involvement, while other modules are automated. The purple modules represent the use of LLMs, where the outputs of LLMs are RDF triples serialized in Turtle. The red modules rely on manually designed prompts.} \label{fig1}
\end{figure}

Here we elaborate the steps of the pipeline for the construction of the soil health KG, as illustrated in Figure~\ref{fig1}.

\subsection{Knowledge sources and conceptual model} \label{sec3_1}
We start with pre-processing the knowledge sources, which are the documents and texts that will be inputted into the LLMs.

The primary knowledge source for our KG is the EU soil health report \citep{doi/10.2800/956606} from the European Environment Agency (EEA). We selected this report because it consolidates current scientific understanding of soil health-related concepts, indicators, thresholds, and assessment methodologies across and beyond EU contexts, and it serves as a reference document for EU soil policy development. Meanwhile, it also provides a visual conceptual model that clearly defines the scope of soil health (see Figure 1.1 of the report). This model is the primary guideline for constructing the upper layer of our soil health KG, ensuring its core structure reflects the established domain understanding. Guided by this conceptual model, we then selected several paragraphs that provide more granular information, such as definitions of soil health-related concepts, their relationships, and quantitative data like thresholds or critical ranges for various indicators. However, the report's textual content does not fully cover all aspects mentioned in its conceptual model; for instance, it explicitly notes limitations in its coverage of soil threats such as wind erosion and salinization. To fill these gaps, we identified an additional EEA report, \textit{The state of soils in Europe} \citep{doi/10.2760/7007291}, as a complementary knowledge source. In total, we selected 105 paragraphs from these two reports and defined one corresponding competency question (CQ) for each, which will be discussed further in Section~\ref{sec3_4}.

We recognize that no single knowledge source, or a limited set of sources, can establish a comprehensive knowledge representation of soil health. Therefore, we also adopted a bottom-up approach to identify missing concepts in the KG. This approach involves using the initial KG as a semantic backbone for a soil knowledge metadata repository through keyword matching, allowing us to progressively identify and incorporate concepts that exist in the broader literature but may be absent from our knowledge sources. This approach will be discussed in more detail in Section~\ref{sec6_1}.

\subsection{The ontology} \label{sec3_2}
The ontology for our KG emerged from a bottom-up construction process guided by the NeOn Methodology \citep{doi:10.3233/AO-150145}. Our approach primarily involved reusing and merging existing ontological resources, which corresponds to Scenario 3 (Reusing ontological resources) and Scenario 5 (Reusing and merging ontological resources) of the NeOn. As shown in Figure~\ref{fig1}, after LLMs generated candidate RDF triples, we first validated the completeness of the extracted entities and relationships. We then searched for standardized classes and properties from established ontologies to assign them.

A variety of ontologies and schemas are in our scope. At a foundational level, we employ the Simple Knowledge Organization System (SKOS) Core \citep{10.5555/1383465.1383467} to structure the conceptual backbone of our KG. SKOS provides the framework for organizing soil-related terms (\texttt{skos:Concept}), their definitions (\texttt{skos:definition}), and the hierarchical relationships between them (\texttt{skos:broader} and \texttt{skos:narrower}). Complementing this, Dublin Core metadata terms \citep{https://doi.org/10.1002/bult.70} are utilized to reference and describe bibliographic resources (\texttt{dct:BibliographicResource}) and other entities (e.g., \texttt{dct:Policy} and \texttt{dct:Standard}) cited in the EU soil health report.

The bulk of the domain entities are represented as instances of \texttt{skos:Concept}, which makes the resulting resource resemble a semantically rich SKOS vocabulary more than a traditional, class-centric OWL ontology. This is a deliberate design choice: our primary use cases (metadata annotation and knowledge discovery, discussed in Section~\ref{sec6_1}) require broad coverage and flexible linking of concepts, rather than a highly axiomatized ontology with strict logical constraints. To bridge flexibility with semantic expressiveness and express the rich and diverse relationships inherent in soil science, we provide a large set of soil-health-specific properties. This allows us to go beyond the core SKOS relations and encode hundreds of domain-relevant predicates between \texttt{skos:Concept} instances, while keeping the overall resource lightweight and adaptable as the KG evolves.

For these more advanced relational needs, we draw upon terms from several ontologies, including Agrontology \citep{SUBIRATSCOLL2022105965}, Semanticscience Integrated Ontology \citep{Dumontier2014}, Open Biological and Biomedical Ontology \citep{Smith2007} (especially Relation Ontology \citep{Smith2005}), Schema.org terms \citep{10.1145/2844544}, SWEET Ontologies \citep{RASKIN20051119}, Wikidata \citep{10.1145/2629489}, Biolink \citep{https://doi.org/10.1111/cts.13302}, BioAssay Ontology \citep{Visser2011-up}, Allotrope Foundation Ontology \citep{MILLECAM20211922}, Time Ontology \citep{10.1145/1017068.1017073}, PROV Ontology \citep{lebo2013prov}, QUDT (Quantities, Units, Dimensions, and Types) \citep{fairsharing_qudt_2025}, and OM (Ontology of Units of Measure) \citep{doi:10.3233/SW-2012-0069}. To systematically manage and apply the properties from these varied ontologies, we categorize the relationships they express into seven main types:

\begin{itemize}
    \item Process relationships: These describe how entities are used within, make use of, or are involved in specific processes (e.g., \\\texttt{agrontology:isUsedIn} and \texttt{obo:RO\_0002331} ``involved in'').
    \item Causative relationships: These capture causal links, such as influences, regulatory effects (positive or negative), correlations, and dependencies between entities (e.g., \texttt{agrontology:causes}, \texttt{obo:RO\_0002213} ``positively regulates'', \texttt{schema:isBasedOn} and \texttt{wdt:P1889} ``different from'').
    \item Quantitative relationships: These are used to link entities to their measured values, including indicators, thresholds and critical intervals (e.g., \texttt{sorelm:hasThreshold} and \texttt{bao:BAO\_0095002} ``has indicator''). QUDT and OM are used to express units of those indicators and quantitative values of their thresholds and ranges (e.g., \texttt{qudt:numericValue} and \texttt{om:hasFactor}).
    \item Partitive relationships: These express part-whole relations such as composition, inclusion and component structures (e.g., \\\texttt{agrontology:isPartOf} and \texttt{agrontology:includes}).
    \item Taxonomic relationships: These are used to classify entities or link them to their different forms (e.g., \texttt{af-x:AFX\_0002806} ``classifies'' and \texttt{agrontology:hasAbbreviation}).
    \item Spatial relationships: These define spatial containment between entities (e.g., \texttt{sio:SIO\_000202} ``contains'') and record geographical context of the data (e.g., the location where thresholds apply using \texttt{prov:atLocation}).
    \item Temporal relationships: These specify the sequence or order of processes in time (e.g., \texttt{agrontology:follows}) and model sampling intervals or observation periods (e.g., \texttt{time:hasDuration}).
\end{itemize}

\subsubsection{Ontology compliance}
Although we strive to reuse established ontologies wherever possible, existing standards alone cannot fully express all relationships specific to our soil health KG. In particular, we adopt Scenario 8 (Restructuring ontological resources) of the NeOn Methodology to define additional properties, when none of the aforementioned ontologies adequately represents our domain needs. For instance, we define \texttt{she:alters}\footnote{We prefix the newly defined entities and relationships in the KG with the \texttt{she} prefix.} and its inverse property \texttt{she:isAlteredBy} for generic usage. Other examples of newly-defined, domain-specific relationships include \texttt{she:sequesters} (describing carbon sequestration processes in soils), \texttt{she:hasEquilibriumWith} (describing chemical balance), \texttt{she:hasDensity} and \texttt{she:hasBiomass} (describing the density and biomass of soil organisms).

By creating such an ontological extension that incorporates these new properties, we ensure that every entity and relationship in the KG is grounded in a formal definition, in accordance with the ontology-compliant approach \citep{10.1007/978-3-031-43458-7_48}. In other words, all classes and properties in the ontology are utilized in our KG.

\subsubsection{Relationship mapping from natural language to ontology properties} \label{sec3_2_2}
A key aspect of our ontology construction process is mapping natural language expressions of relationships to standardized ontology properties. Because human-authored texts utilize varied phrasing for identical conceptual relationships, semantic normalization is crucial \citep{carta2023iterativezeroshotllmprompting}. For example, when modeling the sentence,

\textit{Soil degradation can be defined as a decline in soil quality (Bone et al., 2010), resulting in the reduced functioning of the soil,}

\noindent we use \texttt{obo:RO\_0002212} ``negatively regulates'' to represent the relationship between soil degradation and soil functioning. Regardless of whether the text uses terms like ``reduce'', ``decrease'', or ``decline'', the fundamental relationship is captured by this ontology property. Similarly, verbs like ``cause'', ``lead to'', and ``result in'' are mapped to \texttt{agrontology:causes}. By consolidating these varied lexical forms into a finite set of ontology properties, we preserve semantic clarity and avoid the explosion of bespoke predicates. The complete relationship mapping is shown in Table~\ref{tab1}.

\begin{table}[ht]
\centering
\caption{Relationship mapping from natural language to ontology properties} \label{tab1}
\begin{tabular}{|l|l|}
\hline
\textbf{\begin{tabular}[c]{@{}l@{}}Relationships in natural language\end{tabular}} &
\textbf{\begin{tabular}[c]{@{}l@{}}Ontology property\end{tabular}} \\ \hline

causes, leads to, results in &
\texttt{agrontology:causes} \\
\hline
reduces, decreases, declines &
\begin{tabular}[c]{@{}l@{}}\texttt{obo:RO\_0002212}\\ ``negatively regulates''\end{tabular} \\
\hline
depends on, relies on &
\begin{tabular}[c]{@{}l@{}}\texttt{obo:RO\_0002502}\\ ``depends on''\end{tabular} \\
\hline
assesses, estimates &
\texttt{schema:assesses} \\
\hline
via, through &
\texttt{sorelch:hasMedium} \\
\hline
affects, influences, impacts on &
\texttt{agrontology:affects} \\
\hline
harms, impairs, damages &
\begin{tabular}[c]{@{}l@{}}\texttt{obo:RO\_0002552}\\ ``results in ending of''\end{tabular} \\
\hline
changes, alters &
\texttt{she:alters} \\
\hline
is caused by, due to &
\texttt{agrontology:isCausedBy} \\
\hline
supports, underpins &
\texttt{schema:inSupportOf} \\
\hline
\end{tabular}
\end{table}

Finally, after integrating all classes and properties to form the complete ontology for our KG, we validated its quality and consistency. We used the OOPS! (OntOlogy Pitfall Scanner!) web service \citep{poveda2014oops} to automatically detect common modeling pitfalls. The analysis confirmed that our final ontology is free of detectable pitfalls.

\subsection{Knowledge graph construction}\label{sec3_3}
\subsubsection{Knowledge graph generation using LLMs}
In order to accelerate the construction of our soil health KG, we leverage the LLM (in our case, GPT-5 \citep{2025openaigpt5card}) to generate an RDF representation for every selected paragraph from the knowledge sources. In particular, the process begins with the generation of candidate RDF triples in Turtle syntax from one paragraph of text via prompting. We invoke the model with only the text and the KG generation prompt. After the initial output is generated, we inspect whether any triple retains text as a literal under \texttt{skos:definition} or \texttt{skos:note}. If so, these triples are fed back to the LLM for a secondary extraction, allowing entities and relationships embedded in the retained text to be converted into additional RDF triples, in case they failed to be extracted exhaustively in the first round. The resulting outputs are then parsed by \texttt{rdflib}\footnote{https://github.com/RDFLib/rdflib} to check syntactic correctness. If errors are detected, a Turtle repair mechanism is initiated by feeding the erroneous triples back to the LLM to generate a corrected response, which is then re-validated until there are no syntax errors. The system prompt for candidate RDF triple generation is listed in~\ref{app1}.

\textbf{Blank nodes.} \label{sec3_3_1} In our KG construction, blank nodes play a critical role in representing complex statements that cannot be adequately captured by simple subject-predicate-object triples. We employ blank nodes specifically when one sentence contains more than three elements, a common occurrence in the scientific literature where contextual conditions are frequently expressed. For example, when processing the sentence:

\textit{Soil nutrient status affects biomass production in natural soils and crop yields in agricultural soils,}

\noindent a simple triple structure would fail to preserve the contextual distinction between natural soils and agricultural soils. To maintain this semantics, we use RDF triples with blank nodes shown in Listing~\ref{lst1_1}, in which ``\texttt{obo:BFO\_0000066}'' represents ``occurs in''.

\begin{lstlisting}[style=TurtleStyle, caption={Blank nodes for expressing conditions.}, label={lst1_1}]
she:SoilNutrientStatus agrontology:affects
                [ obo:BFO_0000066 she:NaturalSoils ;
                  rdf:object she:BiomassProduction ],
                [ obo:BFO_0000066 she:AgriculturalSoils ;
                  rdf:object she:CropYields ] .
\end{lstlisting}

Blank nodes also prove essential for expressing thresholds and critical ranges of indicators, another common pattern in soil science literature. For instance, when processing the text about soil organic carbon (SOC):

\textit{If SOC content is 1.2-1.5\%, stability declines rapidly. Lopes et al. (2013) delineated a critical range of SOC for two soil orders: alfisols with 0.5-0.77\% SOC; and entisols with 1.03-1.16\% SOC},

\noindent the modeling results are shown in Listing~\ref{lst1_2}.
\begin{lstlisting}[style=TurtleStyle, caption={Blank nodes for expressing ranges of indicators.}, label={lst1_2}]
she:SOCCriticalRanges af-x:AFX_0002696
                [ agrontology:isUsedIn she:Alfisols ;
                  biolink:has_percentage "0.5-0.77%" ], 
                [ agrontology:isUsedIn she:Entisols ;
                  biolink:has_percentage "1.03-1.16%" ], 
                [ obo:RO_0002212 she:SoilStability ;
                  biolink:has_percentage "1.2-1.5%" ] .
\end{lstlisting}

\noindent Here, ``\texttt{af-x:AFX\_0002696}'' represents ``has range''. This structure preserves the context that different soil types have different critical SOC ranges, and that a specific range declines soil stability.

\subsubsection{Knowledge graph post-processing}
Once syntactically correct, the KG segments undergo several post-processing steps. The first step is to review the generated triples to ensure that all entities and their relationships are correctly and completely extracted from the input text. After that, these candidate triples are manually aligned with ontologies, which we elaborated in Section~\ref{sec3_2}. Then, we integrate these individual aligned KG segments with the RDF representation of the conceptual model mentioned in Section~\ref{sec3_1}. This results in a raw, unified graph that combines the high-level domain structure with detailed, extracted knowledge. This raw graph, however, typically contains duplicate entities and potentially conflicting relationships. Therefore, we then perform entity normalization and relation disambiguation to refine the KG.

\textbf{Entity normalization.} After integrating, each entity within the graph is required to have one preferred label (\texttt{skos:prefLabel}), and optionally one or more alternative labels (\texttt{skos:altLabel}). For every label, we used OpenAI's text-embedding-3-large model \citep{neelakantan2022textcodeembeddingscontrastive} to generate corresponding text embeddings, which serve as high-dimensional vector representations that encode the semantic content of the labels. Subsequently, pairwise comparisons are performed between all entities by calculating the cosine similarity between each pair of label embeddings. For entities with multiple labels, the maximum cosine similarity score among all label pairs is retained as the final similarity measure for that entity pair.

To generate candidate clusters for expert review, we first identify all entity pairs whose similarity score exceeds a pre-defined threshold. We then cluster these pairs based on shared entities. For example, if the similarity score for the pair $(A, B)$ and the pair $(B, C)$ are both above the threshold, then entities $A$, $B$, and $C$ are grouped into a single candidate cluster. Recognizing that the final decision to merge entities must be made by domain experts, our goal was not to find a single optimal threshold, but to use this technique to generate clusters that revealed different patterns of duplicates.

Finally, we collaborated with soil scientists to review these candidate clusters and decide which label to retain as the preferred label, with all others as alternative labels. In this collaboration, we established four main principles for selecting preferred labels. First, we prioritize plural forms over singular forms. Second, we favor terms that explicitly include ``soil'' over those that omit it, for example, selecting ``soil pH'' rather than ``pH'', to maintain domain specificity. Third, we avoid abbreviations in preferred labels, ensuring that ``soil organic carbon contents'' is selected over ``SOC contents'', thereby maximizing clarity and accessibility for users with different expertise levels. Fourth, we prefer the adjectival form over possessive constructions, selecting ``structural soil stability'' rather than ``structural stability of soils'' for more concise representation.

\textbf{Relation disambiguation.} After normalizing entities, we addressed conflicting relationships between entity pairs. For any pair of entities, our integration process may have included multiple conflicting or redundant relationships from different paragraphs. To resolve these conflicts, we implemented the relation disambiguation.

First, we focus on SKOS hierarchical relationships. We address inconsistencies that emerge when integrating multiple KG segments with overlapping concept hierarchies. We identify two primary potential inconsistency patterns. The first involves transitive inconsistencies, which manifest as a clash between associative links and hierarchical links, such as the example in Listing~\ref{lst1_3}. The second concerns multiple broader concept issues, where concepts have more than one parent concept.
To address this, our algorithm constructs a directed graph representation of SKOS broader and narrower relationships and employs graph traversal techniques to identify these patterns.

\begin{lstlisting}[style=TurtleStyle, caption={Example of inconsistent hierarchical relationships.}, label={lst1_3}]
she:Soil  skos:narrower  she:SoilFunctions ;
          skos:related   she:SOCContents .
she:SoilFunctions  skos:narrower  she:SOCContents .
\end{lstlisting}

Then, for other non-hierarchical relationships, we focus on identifying and resolving two potential forms of relationship inconsistency across URI-to-URI connections (i.e., ignoring blank nodes and literals). First, we detect multiple relationship instances where different predicates connect the same subject-object pairs, which indicate potential semantic redundancy. Second, we identify bidirectional applications of the same predicate between concept pairs, excluding naturally symmetric predicates (e.g., \texttt{skos:related}). This approach recognizes that while some predicates are inherently symmetric, others require distinct inverse predicates.

Similar to our approach for entity normalization, the goal of relation disambiguation is not to automatically resolve all conflicts but to recognize common conflict patterns and reveal potential inconsistencies for expert review.

\subsubsection{Knowledge graph enrichment} \label{sec3_3_3}
Finally, we enrich the KG by materializing reversible relationships. If the predicate of a triple (subject-predicate-object) is invertible, we add the inverse triple (object-inversePredicate-subject) to the KG to ensure completeness. In addition, the KG is further enriched by interlinking to the controlled vocabularies and thesauri in the field of soil science, including AGROVOC \citep{SUBIRATSCOLL2022105965}, GloSIS \citep{doi:10.1177/22104968251363767}, ISO 11074 \citep{ISO11074}, GEMET \footnote{www.eionet.europa.eu} and INRAE Thesaurus \citep{aubin:hal-03824794}, thereby supporting downstream knowledge discovery, which will be discussed in Section~\ref{sec6_1}. This interlinking is achieved using SKOS mapping properties, specifically \texttt{skos:exactMatch} and \texttt{skos:closeMatch}. We employ \texttt{skos:exactMatch} when the concept's \texttt{skos:prefLabel} in our KG exactly matches the concept's \texttt{skos:prefLabel} in the target vocabulary. In other cases, such as when a \texttt{skos:altLabel} in our KG matches a \texttt{skos:prefLabel} in the vocabulary (or vice versa), \texttt{skos:closeMatch} is used.


\subsection{Human-in-the-loop validation} \label{sec3_4}
To validate our soil health KG, we tested its ability to answer a set of CQs that were defined a priori to guide the construction of the KG itself, as shown in Figure~\ref{fig1}. If the KG faithfully models the knowledge from the source text, it should provide answers consistent with those expected from the original document. Below, we detail our two-step validation process.

\subsubsection{Competency question formulation}
The CQs were defined before the construction of the KG and its ontology to guide the knowledge modeling process. For each of the 105 paragraphs selected from the knowledge sources to be fed into our pipeline, we concurrently formulated a corresponding CQ designed to capture the key information within that text segment. In total, 105 CQs were defined to reflect soil scientists' interests and the scope of soil health as defined by the conceptual model. Crucially, these CQs were designed based on the explicit information present in the knowledge source, ensuring that their correct answers could be directly found within the source text without requiring external information or complex reasoning. Consequently, the ability of our KG to correctly retrieve and provide the relevant information for these CQs serves as a direct measure of how accurately it models the knowledge contained within the source document. All CQs were reviewed by soil scientists to ensure that each question is valid, unambiguous, and relevant within the context of soil science. More examples of CQs are provided in Table~\ref{tab2}.

\subsubsection{SPARQL query translation and answer validation}
Once the set of CQs is established, each question is translated into a SPARQL query designed to retrieve the relevant information from the soil health KG. We then execute these SPARQL queries. The direct results returned by the KG often consist of specific facts or data fragments. Therefore, these retrieved pieces of information are compiled and rewritten into a human-readable answer that directly addresses the original CQ. It is important to note that no additional knowledge is introduced during this answer formulation process; only the information explicitly returned by the queries is used. These formulated answers, along with the original CQs and the KG's query results, are then compared with the information found directly in the knowledge sources to validate both coverage and accuracy. Finally, soil scientists review these answers to confirm whether they hold true from a scientific perspective, helping to catch any misinterpretations of the KG's output.

Here is an example of our validation processes. First, we have a CQ:

\textit{What is the difference between soil health and soil quality?}

\noindent Then we translate this CQ into SPARQL query, as shown in Listing~\ref{lst1_4}, and execute it.

\begin{lstlisting}[style=SPARQLStyle, caption={SPARQL query for the difference between soil health and soil quality.}, label={lst1_4}]
SELECT DISTINCT ?concept ?uniqueDefinition
WHERE {
  ?concept skos:prefLabel ?label ;
           skos:definition ?uniqueDefinition .
  FILTER(?label = "soil health"@en || ?label = "soil quality"@en)
  
  FILTER NOT EXISTS {
    ?otherConcept skos:prefLabel ?otherLabel ;
                  skos:definition ?uniqueDefinition .
    FILTER(?otherLabel != ?label && 
           (?otherLabel = "soil health"@en || ?otherLabel = "soil quality"@en))
  }
}
\end{lstlisting}

\noindent The results returned by the KG are shown in Listing~\ref{lst1_5}. They have been verified and are consistent with the expected answer.

\begin{lstlisting}[style=TurtleStyle, caption={Query result.}, label={lst1_5}]
she:SoilHealth: Soil health is defined as the actual capacity of a soil to perform its core functions and to provide ecosystem services.
she:SoilQuality: Soil quality is defined as the inherent capacity which provides the basic frame within which a soil evolves.
\end{lstlisting}

Following these iterative validation processes, we performed the answer alignment step, as illustrated in Figure~\ref{fig1}. A small number of the triples in the KG were adjusted to ensure that the information returned by the KG for each CQ is consistent with the expected answers, accurately reflecting the corresponding content of the knowledge sources.

\begin{table}[ht!]
\centering
\caption{Examples of competency questions. The CQs are categorized into four types, focusing on the fundamental concepts of soil health, function and properties, soil threats and degradation, their general and specific indicators, and the thresholds and critical ranges of the indicators.} \label{tab2}
\resizebox{\linewidth}{!}{
\begin{tabular}{|l|l|}
\hline
\multirow{2}{*}{\textbf{Question type}} &
  \multirow{2}{*}{\textbf{Example CQs}} \\
 &
   \\ \hline
\multirow{5}{*}{\begin{tabular}[c]{@{}l@{}}Fundamental\\concepts of\\soil health,\\functions,\\ \& properties\end{tabular}} &
  \begin{tabular}[c]{@{}l@{}}What are the components of soil, and what are the\\ relationships between soil and its functions?\end{tabular} \\ \cline{2-2} 
 &
  \begin{tabular}[c]{@{}l@{}}What are soil functions and what is the relationship between\\ soil functions, soil processes, and ecosystem services?\end{tabular} \\ \cline{2-2} 
 &
  \begin{tabular}[c]{@{}l@{}}What are the two types of soil nutrients and\\ what substances do they contain?\end{tabular} \\ \hline
\multirow{6}{*}{\begin{tabular}[c]{@{}l@{}}Soil threats \&\\ degradation\end{tabular}} &
  \begin{tabular}[c]{@{}l@{}}What are soil threats and what factors indicate the existence\\ of soil threats?\end{tabular} \\ \cline{2-2} 
 &
  \begin{tabular}[c]{@{}l@{}}What are the main types of soil degradation and where does\\ soil sealing fall into this category?\end{tabular} \\ \cline{2-2} 
 &
  \begin{tabular}[c]{@{}l@{}}What are the two types of soil pollution, what are the\\characteristics of diffuse soil pollution, \\and what are its typical origins?\end{tabular} \\ \hline
\multirow{5}{*}{\begin{tabular}[c]{@{}l@{}}General \&\\specific\\ indicators\end{tabular}} &
  \begin{tabular}[c]{@{}l@{}}How are indicators defined for assessing nitrogen and phosphorus\\ status in agricultural soils, and what specific types of indicators\\ are mentioned for nitrogen assessment in forest soils?\end{tabular} \\ \cline{2-2} 
 &
  \begin{tabular}[c]{@{}l@{}}What is soil erosion, what is identified as the main\\ indicator used to quantify soil erosion and what is its unit?\end{tabular} \\ \cline{2-2} 
 &
  \begin{tabular}[c]{@{}l@{}}What are the main indicators that describe topsoil\\ compaction and what are their symbols, if any?\end{tabular} \\ \hline
\multirow{3}{*}{\begin{tabular}[c]{@{}l@{}}Thresholds \&\\ critical ranges\end{tabular}} &
  \begin{tabular}[c]{@{}l@{}}What are the critical levels or thresholds of SOC for maintaining\\soil health, structural stability, and optimal crop yield,\\ and how do these thresholds vary across different soil types?\end{tabular} \\ \cline{2-2} 
 &
  \begin{tabular}[c]{@{}l@{}}What is the soil pH range that can constrain crop\\ production and what are the critical and optimal pH value?\end{tabular} \\ \hline
\end{tabular}}
\end{table}

\section{Results}\label{sec4}
\subsection{Overview of the soil health KG}\label{sec4_1}
\begin{figure}
\includegraphics[width=\textwidth]{soil_health_KG.pdf}
\caption{Overview of the high-level structure of the soil health KG, colors represent concepts at different hierarchy levels in the KG. The double-rounded rectangle indicates that the concept exists in other soil-related thesauri. Note that all concepts and their relationships have been omitted from the prefixes and replaced with their human-readable labels for better visualization.} \label{fig2}
\end{figure}

Figure~\ref{fig2} illustrates the high-level structure of our KG, which is derived from the conceptual model. The resulting soil health KG comprises a total of 11,719 RDF triples, which describe 2,017 distinct entities. The majority of these entities (1,785) represent soil-related concepts, 493 of which are also matched in established thesauri. The remaining entities include 158 bibliographic references cited in the knowledge sources, along with other supportive entities such as relevant policies, standards, regulations, geographical regions (e.g., cities, continents), and mathematical or quantitative entities (including models, equations, and units of measure). Structurally, the KG leverages 19 classes and 205 properties drawn from the selected ontologies, of which 45 properties are defined by us. The resource availability statement is provided at the end of this paper.

\subsection{Knowledge graph generation results} \label{sec4_2}
To evaluate the quality of LLM-generated RDF triples, we compared the generated graphs against human-labeled gold standard graphs across all 105 paragraphs. The LLM-generated triples are produced by the system prompt in~\ref{app1}, using GPT-5 (default settings) with a single attempt per paragraph (i.e., no self-consistency voting \citep{wang2023selfconsistency}). For each paragraph, we define the LLM-generated RDF triples as a generated graph $G_{gen}$ and the human-annotated RDF triples as a gold standard graph $G_{gold}$. We employed a suite of complementary metrics to assess different aspects of generation quality at the triple level \citep{ghanem:hal-04862235}.

For evaluation purposes, we convert RDF triples into directed graphs using NetworkX \texttt{MultiDiGraph} structures \citep{osti_960616}. In this representation, subjects and objects serve as nodes, while predicates form directed edges from subjects to objects. The \texttt{MultiDiGraph} structure accommodates parallel edges, allowing multiple predicates to connect the same subject-object pair, a common occurrence in RDF where entities may have multiple relationships. For n-gram based metrics, each triple $(s, p, o)$ is linearized into a token sequence by concatenating the subject, predicate, and object labels, enabling the application of text-based evaluation methods to graph structures.

\textbf{Exact match metrics.} Triple match Precision, Recall, and $F1$ score measure the extent to which generated triples exactly match gold standard triples. These metrics provide a strict evaluation of the LLM's ability to reproduce the precise subject-predicate-object structures defined by human annotators.

\textbf{N-gram overlap metrics.} G-ROUGE and G-BLEU adapt traditional text generation metrics to the graph domain by treating linearized triples as sequences. G-ROUGE (using ROUGE-2, bi-gram) captures recall-oriented overlap, measuring how well the generated graph covers the content of the gold standard. G-BLEU (using 4-gram) provides a precision-oriented perspective, evaluating the fluency and accuracy of generated triple sequences. These metrics offer a more lenient evaluation than exact matching by crediting partial overlaps in entity and relation naming.

\textbf{Semantic similarity metric.} Graph BERTScore leverages contextual embeddings from a pre-trained language model (in our case, RoBERTa-large \citep{liu2019robertarobustlyoptimizedbert}) to compute semantic similarity between triples. This metric captures cases where generated triples convey equivalent meaning despite lexical differences, such as when synonymous terms or paraphrased relations are used.

\textbf{Structural metric.} Graph Edit Distance (GED) quantifies the minimum number of edit operations (node insertions, deletions, and substitutions; edge insertions, deletions, and substitutions) required to transform the generated graph into the gold standard graph. This metric evaluates the structural fidelity of the generated graph to the gold standard graph.

The complete mathematical definitions of the aforementioned metrics are provided in~\ref{app2}. Table~\ref{tab3} presents the aggregate evaluation results across all 105 paragraph-graph pairs.

\begin{table}[ht]
\centering
\caption{Evaluation results of LLM-generated knowledge graphs} \label{tab3}
\begin{tabular}{|l|c|c|c|}
\hline
\textbf{Metric} & \textbf{Precision} & \textbf{Recall} & \textbf{$F1$} \\ \hline
Triple Match & 0.2465 & 0.2548 & 0.2506 \\ \hline
G-Rouge (ROUGE-2) & 0.6478 & 0.6703 & 0.6588 \\ \hline
G-BLEU (4-gram) & 0.6913 & 0.7153 & 0.7031 \\ \hline
Graph BERTScore (RoBERTa-large) & 0.8456 & 0.8749 & 0.8600 \\ \hline
\textbf{Metric} & \multicolumn{3}{c|}{\textbf{Mean}} \\ \hline
Graph Edit Distance (GED) & \multicolumn{3}{c|}{0.4730} \\ \hline
\end{tabular}
\end{table}


\subsection{Knowledge graph validation results}\label{sec4_3}
We also conducted a quantitative evaluation of the semantic alignment between the KG outputs and expected answers. We executed all 105 SPARQL queries against the final KG and compared their outputs with the expected answers using two complementary metrics: BERTScore \citep{Zhang*2020BERTScore:} for token-level similarity and semantic textual similarity (STS) \citep{cer-etal-2017-semeval} for sentence-level similarity using OpenAI's text-embedding-3-large model. Table~\ref{tab4} presents the evaluation results.

\begin{table}[ht]
\centering
\caption{Statistical results of the alignment between KG outputs and expected answers} \label{tab4}
\begin{tabular}{|l|ccc|}
\hline
\textbf{Metric}    & \multicolumn{1}{c|}{\textbf{Precision}} & \multicolumn{1}{c|}{\textbf{Recall}} & \textbf{$F1$} \\ \hline
BERTScore & \multicolumn{1}{c|}{0.8088}             & \multicolumn{1}{c|}{0.7779}          & 0.7923            \\ \hline
\textbf{Metric} & \multicolumn{3}{c|}{\textbf{Average similarity}} \\ \hline
Semantic textual similarity    & \multicolumn{3}{c|}{0.8945}                      \\ \hline
\end{tabular}
\end{table}


\section{Discussion}\label{sec5}
\subsection{Evaluation of LLM-generated knowledge graphs}\label{sec5_1}
The results in Table~\ref{tab3} show a clear contrast between exact structural agreement and semantic similarity. The low Triple Match $F1$ indicates that LLM-generated triples rarely reproduce the gold-standard graphs exactly. In contrast, the higher G-ROUGE, G-BLEU, and Graph BERTScore values show that the generated graphs often preserve much of the relevant content and meaning of the source text. Taken together, these results suggest that the LLM is generally effective at extracting domain knowledge, but less consistent in expressing that knowledge in the same graph structure as human annotators.

The mean normalized GED of 0.4730 provides complementary insight into the structural differences between the generated and gold standard graphs. While the semantic content aligns well, this GED value indicates a substantial degree of structural divergence, requiring edit operations to transform the generated graph into the gold standard. In practical terms, human annotators can start from LLM-generated drafts that already capture semantically similar content, then focus on structural corrections and ontology-aligned reformulation, rather than constructing the KG entirely from scratch. This reduction in manual effort demonstrates the value of LLMs as an accelerator for knowledge engineering tasks, even when their outputs require subsequent refinement.


Our analysis reveals that much of this structural divergence arises from representation choices rather than semantic failure. In particular, we observed two recurring patterns. First, regarding entity granularity, the LLM tends to merge related concepts into single entities, whereas human annotators ensure that each URI represents only one distinct concept. For example, the LLM usually represents ``inputs of nutrients and contaminants'' as a single node, while human annotation would maintain ``inputs of nutrients'' and ``inputs of contaminants'' as separate entities. Second, regarding variations in knowledge modeling, we observe that the LLM tends to overuse blank nodes for expressing conditional statements, even when direct subject-predicate-object triples would suffice. As described in Section~\ref{sec3_3_1}, blank nodes are employed in our pipeline to represent complex statements with contextual qualifiers; however, the LLM frequently applies this complex structure to straightforward factual statements that do not require such reification. These patterns collectively result in LLM-generated graphs exhibiting structural differences from human-annotated graphs, despite conveying similar semantic content. These findings reinforce the necessity of human-in-the-loop post-processing in our pipeline, where expert review ensures not only semantic accuracy but also appropriate entity granularity and structural alignment with established ontological conventions.

\subsection{Validation of the soil health KG}\label{sec5_2}
The validation results presented in Table~\ref{tab4} demonstrate strong alignment between the KG outputs and expected answers derived from the source documents. The STS metric achieves an average score of nearly 0.9, confirming that the KG outputs and expected answers are strongly aligned semantically. The BERTScore $F1$, with an average close to 0.8, is lower than the STS score. This difference is expected and reasonable: BERTScore operates at the token level, and as discussed in Section~\ref{sec3_2_2}, our methodology includes relationship mapping that transforms varied natural language expressions into standardized ontology properties. While such normalization preserves semantic meaning, as evidenced by the high semantic similarity scores, it naturally introduces token-level differences between the KG representation and the source text. These quantitative results validate that our KG accurately captures and represents the knowledge from the source documents while maintaining semantic fidelity through ontological standardization.

Moreover, the KG, by its nature, structures and aggregates knowledge from different paragraphs within the knowledge sources. Consequently, when executing queries, we frequently observed that the KG provides information drawn from various locations across the knowledge sources, instead of just the expected answers from a single paragraph source. For example, we have another CQ:

\textit{What knowledge is needed in order to manage soil sustainably and protect it?}

\noindent And we executed its corresponding SPARQL query, the result is shown in Listing~\ref{lst1_6}.

\begin{lstlisting}[style=TurtleStyle, caption={Query result, \texttt{obo:T4FS\_0000577} represents ``requires knowledge about''}, label={lst1_6}]
she:SoilProtection obo:T4FS_0000577 she:CriticalLevelAtPreventiveActivity
she:SoilProtection obo:T4FS_0000577 she:CriticalLevelAtRestorativeActivity
she:SoilProtection obo:T4FS_0000577 she:CriticalLevelAtSoilFunctions
she:SoilProtection obo:T4FS_0000577 she:StateOfSoils
she:SustainableSoilManagement obo:T4FS_0000577 she:StateOfSoils
she:StateOfSoils agrontology:isAffectedBy she:ClimaticConditions
she:StateOfSoils agrontology:isAffectedBy she:SoilManagement
\end{lstlisting}
However, the expected answer to this CQ is ``In order to manage soil sustainably and protect it, knowledge is needed about the state of soil and how it develops under current and future management and climatic conditions.'', which originates from a single paragraph.
Clearly, for this CQ, the KG output not only covers all the information in the expected answer but also provides information synthesized from other paragraphs. This further explains why BERTScore is lower than STS, as the KG sometimes outputs more information than the expected answer. Also, this illustrates the KG’s ability to integrate and reorganize distributed information across the source corpus.

\section{Applications of the soil health KG}\label{sec6}
\subsection{Knowledge discovery in soil science} \label{sec6_1}
Once the soil health KG is established, its first and foremost application is to serve as the semantic backbone for a soil knowledge metadata repository.  We, and the SoilWise project\footnote{https://soilwise-he.eu/} as a whole, aim to create a repository by interlinking external data and knowledge to the soil health KG. To achieve this, we harvest a vast array of soil-related metadata records from general platforms such as CORDIS\footnote{https://cordis.europa.eu/}, OpenAire\footnote{https://www.openaire.eu/}, and soil-specific platforms like PREPSOIL\footnote{https://prepsoil.eu/}. This harvested metadata is then serialized into RDF, forming a metadata KG. The primary method to integrate the soil health KG with this metadata KG is keyword matching. In this process, each concept within the soil health KG serves as a keyword. These keywords are then matched against keywords present in the metadata records to establish bi-directional links, thereby interlinking the two KGs.

This combined knowledge repository can immediately be used as an interactive knowledge discovery tool. The SoilWise data and knowledge hub\footnote{https://voc.soilwise-he.containers.wur.nl/} functions as a soil health thesaurus, acting as an entry point to multiple domain thesauri by linking our KG to established resources discussed in Section~\ref{sec3_3_3}. In addition, by integrating with the metadata KG, the portal also displays metadata of peer-reviewed papers, datasets, and other scientific outputs containing these concepts as keywords. This transforms the portal into a convenient gateway for researchers, allowing them to find both standardized vocabularies and their related scientific outputs in a single interface.

Furthermore, such integration also helps us identify concepts that are missing in our soil health KG but are widely used and researched in the soil science community. During the development of the KG, we employed keyword matching in parallel. Therefore, we had a list of unmatched keywords from metadata records. We then collaborate with soil experts to evaluate if some of these keywords should appear in the soil health KG. A prime example is that literature focusing on soil health usually includes limited information on soil classification. Therefore, we utilized another EEA report \citep{doi/10.2760/7007291} to fill this gap.

\subsection{A benchmark for text2KG}\label{sec6_2}
Besides the use case discussed above, our work contributes a dataset mapping text segments from scientific literature to corresponding RDF triples. The construction of KGs from unstructured text remains a challenging task. While existing benchmarks like WebNLG \citep{gardent-etal-2017-webnlg} and T-REx \citep{elsahar-etal-2018-rex} are valuable, they predominantly focus on basic triple extraction from general-purpose sources like DBpedia and Wikipedia. Our dataset, by contrast, is derived from technical, domain-specific literature, which presents unique and complex modeling challenges.

As detailed in Section~\ref{sec3_3_1}, scientific texts frequently express complex information with specific contextual conditions (e.g., ``biomass production in natural soils'' versus ``crop yields in agricultural soils''). To accurately capture this semantics, our dataset makes widespread use of reification with blank nodes. This heavy reliance on blank nodes to represent conditional knowledge, combined with the domain-specific terminology and complex sentence structures inherent in academic literature, makes our dataset a challenging benchmark. It is well-suited for evaluating the robustness of text-to-knowledge-graph (text2KG) methods, particularly those that leverage LLMs, as these models may struggle to correctly and consistently generate such complex, reified RDF structures while maintaining ontological alignment. Table~\ref{tab5} provides basic information about this dataset, including the number of text segments, average words, resulting triples, unique entities and relationships.

\begin{table}[ht]
\centering
\caption{Information of the text2KG benchmark dataset. "\#" represents number.} \label{tab5}
\begin{tabular}{|llllll|}
\hline
\multicolumn{2}{|l|}{\textbf{\# of paragraphs}} & \multicolumn{2}{l|}{\textbf{Average \# of words}} & \multicolumn{2}{l|}{\textbf{Average \# of triples}} \\ \hline
\multicolumn{2}{|l|}{105} & \multicolumn{2}{l|}{140.82} & \multicolumn{2}{l|}{92.08} \\ \hline
\multicolumn{3}{|l|}{\textbf{Average \# of entities}}                     & \multicolumn{3}{l|}{\textbf{Average \# of relationships (properties)}}        \\ \hline
\multicolumn{3}{|l|}{28.62}                & \multicolumn{3}{l|}{15.52}             \\ \hline
\end{tabular}
\end{table}


\section{Conclusion}\label{sec7}
In this paper, we addressed the need for structured and accessible soil health knowledge by presenting two main contributions: an ontology-compliant soil health KG and the semi-automated, human-in-the-loop pipeline developed for its construction. Our work shows that LLMs can support knowledge engineering by accelerating the generation of candidate RDF graphs from scientific text, while expert review remains essential for ontology alignment, structural normalization, and validation.

The evaluation results indicate that LLM-generated graphs often capture the semantic content of the source text well, but still differ substantially from human-annotated graphs in exact and structural representation. This confirms the value of LLMs as accelerators in KG construction, while also highlighting the continuing importance of human refinement in ontology-sensitive domains. The resulting soil health KG was reviewed through competency-question-based validation and supports practical applications in knowledge discovery, while also providing a text-to-RDF resource for future benchmarking.

Future work will focus on two complementary directions. First, we will further improve the pipeline itself by exploring alternative prompting strategies, retrieval-augmented generation settings, and the fine-tuning of smaller language models for ontology-compliant RDF generation. Second, we will investigate the transferability of the proposed workflow to other scientific domains, in order to assess how well it generalizes beyond soil health and which components remain domain-specific. More broadly, we see this work as a step toward more scalable and reliable approaches for semi-automated knowledge engineering.


\section*{Data availability}
All resources described in this paper are hosted and maintained by the SoilWise project to ensure their long-term availability. The soil health KG, its supporting ontology schema and the validation dataset with CQs are available at \url{https://doi.org/10.5281/zenodo.14936019}. The AgroPortal instance is accessible at \url{https://agroportal.lirmm.fr/ontologies/SHKG}. The WIDOCO HTML page can be accessed at \url{https://soilwise-he.github.io/soil-health}, where all self-defined URIs are resolved to. The SPARQL endpoint is provided at \url{https://repository.soilwise-he.eu/sparql/}. Finally, the code and prompts for the KG construction pipeline can be found at \url{https://github.com/soilwise-he/soil-health-knowledge-graph}. All resources related to this project are licensed under the Creative Commons Attribution 4.0 International (CC BY 4.0) license\footnote{https://creativecommons.org/licenses/by/4.0/}.

\section*{Acknowledgments}
This work was supported by the EU’s Horizon Europe research and innovation programme within the SoilWise project (grant agreement ID: 101112838). We thank all soil scientists involved for their invaluable input and domain expertise throughout this work.

%% The Appendices part is started with the command \appendix;
%% appendix sections are then done as normal sections
\appendix
\section{The system prompt for KG generation}
\label{app1}
\begin{markdownprompt}
# RDF Knowledge Graph Instructions for Constructing Soil Health Knowledge Graph

## 1. Overview
You are a top-tier algorithm designed for extracting information from unstructured text to build an RDF knowledge graph in the domain of **soil science and soil health**.
Your task is to identify entities, relationships and references from a given text, which will be a **single paragraph from a soil health report**.
Try to capture as much information from the text as possible without sacrificing accuracy. Do not add any information that is not explicitly mentioned in the text.

## 2. Subjects
- **Subjects** represent entities and concepts relevant to **soil science and soil health**.
- Keep subjects atomic and distinct; split combined mentions into separate subjects and do not include conjunctions such as "and" inside a single subject.
- Attempt to extract as many entities and concepts as you can within this domain.
- The aim is to achieve simplicity and clarity in the knowledge graph, making it accessible for a vast audience.

## 3. Labeling Subjects
- **Consistency**: Ensure you assign at least one available class from existing ontologies to each subject. **Do not customize classes**.
  - When you identify a subject as a concept, always use `'skos:Concept'` to label it.
  - If a subject is not a concept, please search for a class from an extant ontology to label it.
  - For every entity and concept, include one `skos:prefLabel` with the human-readable name in lowercase found in the text.
- **URIs**: URIs should be names or human-readable identifiers derived from the text. Use `she: <https://soilwise-he.github.io/soil-health#>` as prefix. Split words in URIs by capitalizing initial letters.
   
## 4. Relations and Predicates
- **Relationships** represent connections between entities.
- Attempt to extract as many relations as you can.
- Ensure consistency and generality in relationship predicates when constructing knowledge graphs. Instead of using specific and momentary predicates, **make sure to use general and timeless relationship predicates.**

## 5. Nested Triples with Blank Nodes
- **Blank Nodes**: When a relationship requires **multiple contextual qualifiers** (e.g., different **conditions** or **locations**), use blank nodes to create nested triples.
- Use blank nodes only when a relationship requires multiple properties and a single predicate-object pair cannot fully capture the text's nuance.

## 6. Coreference Resolution
- **Maintain Entity/Concept Consistency**: When extracting entities and concepts, it's vital to ensure consistency.
- If an entity or concept is mentioned multiple times in the text but referred to by different forms, always use the most complete identifier and preferred label for that entity or concept throughout the knowledge graph.

## 7. Objects
- **Retaining Text as Objects**: Do not retain the given text as objects, except in the following case:
  - **Definitions**: If the text is defining concepts, include the full definition as an object and use `'skos:definition'` as the predicate.
- **Important**: Even when you retain parts of the text as objects (definitions), you must still extract entities, relationships, and references according to the previous instructions.

## 8. References and Citations
- **Extract References**: Identify any references or citations that appear in the text.
- **Labeling References**:
  - Use Dublin Core (`dcterms`) to label references.
  - Represent each reference as a `dcterms:BibliographicResource`.
  - Do not include `skos:prefLabel` for any references, instead use properties from Dublin Core.
- **Connecting Information to References**:
  - Use `dcterms:source` to indicate which information in the text references this literature.
  - When an extracted piece of information is supported by a reference, link it to the corresponding `dcterms:BibliographicResource` using `dcterms:source`.

## 9. Strict Compliance
Adhere to the rules strictly. Non-compliance will result in termination.
\end{markdownprompt}

\section{Mathematical definitions of the evaluation metrics} \label{app2}
This appendix provides formal definitions for the evaluation metrics used to assess knowledge graph generation quality in Section~\ref{sec4_2}. Let $G_{gen} = (V_{gen}, E_{gen})$ denote the generated graph and $G_{gold} = (V_{gold}, E_{gold})$ denote the gold standard graph, where $V$ represents the set of nodes (entities) and $E$ represents the set of edges (triples). Each triple $t$ is represented as a tuple $(s, p, o)$ consisting of subject $s$, predicate $p$, and object $o$. We denote the set of triples in the generated graph as $T_{gen}$ and in the gold standard graph as $T_{gold}$.

\subsection{Triple Match metrics}\label{app2_1}
Triple Match metrics evaluate exact correspondence between generated and gold standard triples.

\textbf{Precision} measures the fraction of generated triples that exactly match triples in the gold standard:

$$
P_{match} = \frac{|T_{gen} \cap T_{gold}|}{|T_{gen}|}
$$

\textbf{Recall} measures the fraction of gold standard triples that are exactly reproduced in the generated graph:

$$
R_{match} = \frac{|T_{gen} \cap T_{gold}|}{|T_{gold}|}
$$

\textbf{$F1$ Score} is the harmonic mean of precision and recall:

$$
F1_{match} = \frac{2 \cdot P_{match} \cdot R_{match}}{P_{match} + R_{match}}
$$

\subsection{G-ROUGE (ROUGE-2)}\label{app2_2}
G-ROUGE adapts the ROUGE-2 metric to evaluate bi-gram overlap between linearized graph representations. Each triple $t = (s, p, o)$ is linearized into a token sequence by concatenating the subject, predicate, and object labels. Bi-grams (consecutive two-token sequences) are then extracted from these linearized triples.

Let $N_2(T)$ denote the multiset of bi-grams extracted from the linearized representation of triple set $T$. Precision, Recall, and $F1$ are computed as:

$$
P_{rouge} = \frac{|N_2(T_{gen}) \cap N_2(T_{gold})|}{|N_2(T_{gen})|}
$$

$$
R_{rouge} = \frac{|N_2(T_{gen}) \cap N_2(T_{gold})|}{|N_2(T_{gold})|}
$$

$$
F1_{rouge} = \frac{2 \cdot P_{rouge} \cdot R_{rouge}}{P_{rouge} + R_{rouge}}
$$

where $\cap$ denotes multiset intersection.

\subsection{G-BLEU (4-gram)}\label{app2_3}
G-BLEU adapts the BLEU metric to evaluate 4-gram overlap between linearized graph representations. Each triple is linearized into a token sequence, and 4-grams (consecutive four-token sequences) are extracted.

Let $N_4(T)$ denote the multiset of 4-grams extracted from the linearized representation of triple set $T$. Precision, Recall, and $F1$ are computed as:

$$
P_{bleu} = \frac{|N_4(T_{gen}) \cap N_4(T_{gold})|}{|N_4(T_{gen})|}
$$

$$
R_{bleu} = \frac{|N_4(T_{gen}) \cap N_4(T_{gold})|}{|N_4(T_{gold})|}
$$

$$
F1_{bleu} = \frac{2 \cdot P_{bleu} \cdot R_{bleu}}{P_{bleu} + R_{bleu}}
$$

where $\cap$ denotes multiset intersection.

\subsection{Graph BERTScore}\label{app2_4}
Graph BERTScore measures semantic similarity between triples using contextual embeddings. Let $\phi(t)$ denote the embedding of triple $t$ obtained from a pre-trained language model (RoBERTa-large in our experiments).

For each triple $t_i \in T_{gen}$, we compute its maximum cosine similarity to any triple in $T_{gold}$:

$$
sim(t_i, T_{gold}) = \max_{t_j \in T_{gold}} \frac{\phi(t_i) \cdot \phi(t_j)}{|\phi(t_i)| |\phi(t_j)|}
$$

\textbf{Precision} averages the maximum similarity scores for generated triples:

$$
P_{BERT} = \frac{1}{|T_{gen}|} \sum_{t_i \in T_{gen}} sim(t_i, T_{gold})
$$

\textbf{Recall} averages the maximum similarity scores for gold standard triples:

$$
R_{BERT} = \frac{1}{|T_{gold}|} \sum_{t_j \in T_{gold}} sim(t_j, T_{gen})
$$

where $sim(t_j, T_{gen}) = \max_{t_i \in T_{gen}} \frac{\phi(t_j) \cdot \phi(t_i)}{|\phi(t_j)| |\phi(t_i)|}$

\textbf{$F1$ Score} is the harmonic mean:

$$
F1_{BERT} = \frac{2 \cdot P_{BERT} \cdot R_{BERT}}{P_{BERT} + R_{BERT}}
$$

\subsection{Graph Edit Distance (GED)}\label{app2_5}
Graph Edit Distance quantifies the minimum cost of transforming $G_{gen}$ into $G_{gold}$ through a sequence of edit operations. The set of allowable operations includes:

\begin{itemize}
    \item Node insertion: $ins_v$
    \item Node deletion: $del_v$
    \item Node substitution: $sub_v$
    \item Edge insertion: $ins_e$
    \item Edge deletion: $del_e$
    \item Edge substitution: $sub_e$
\end{itemize}

Let $\mathcal{P}(G_{gen}, G_{gold})$ denote the set of all valid edit paths transforming $G_{gen}$ into $G_{gold}$. For an edit path $p = (o_1, o_2, \ldots, o_k)$, the cost is:

$$
cost(p) = \sum_{i=1}^{k} c(o_i)
$$

where $c(o_i)$ is the cost of operation $o_i$.

The Graph Edit Distance is defined as:

$$
GED(G_{gen}, G_{gold}) = \min_{p \in \mathcal{P}(G_{gen}, G_{gold})} cost(p)
$$

To enable comparison across graphs of different sizes, we compute a normalized GED for each paragraph:

$$
GED_{norm}(G_{gen}, G_{gold}) = \frac{GED(G_{gen}, G_{gold})}{|V_{gen}| + |V_{gold}| + |E_{gen}| + |E_{gold}|}
$$

The overall GED score is computed as the arithmetic mean of normalized GED across all $N$ paragraphs:

$$
\overline{GED} = \frac{1}{N} \sum_{i=1}^{N} GED_{norm}(G_{gen}^{(i)}, G_{gold}^{(i)})
$$

where $G_{gen}^{(i)}$ and $G_{gold}^{(i)}$ denote the generated and gold standard graphs for the $i$-th paragraph, respectively.


%% If you have bib database file and want bibtex to generate the
%% bibitems, please use
%%
\bibliographystyle{elsarticle-num-names} 
\bibliography{reference.bib}

%% else use the following coding to input the bibitems directly in the
%% TeX file.

%% Refer following link for more details about bibliography and citations.
%% https://en.wikibooks.org/wiki/LaTeX/Bibliography_Management

\end{document}

\endinput
%%
%% End of file `elsarticle-template-num-names.tex'.
